\section{The Opposite of a Misconception}

\begin{quotex}
Destroy a man's illusions and you destroy his happiness \flright{\textsc{Hubert van Zeller}}

\end{quotex}
\paragraph{Metaphysical Errors}
I recently stumbled upon an article on an obscure blog titled ``Common Metaphysical Errors", by a writer apparently unaware of the oxymoron involved with that expression. Rather, it is a collection of carefully selected opinions gleaned from unrevealed sources followed with a pithy counterpoint.

Of course, there can be no metaphysical errors. Metaphysics is not a collection of beliefs or opinions, but rather it derives from exacting and precise intellectual intuitions. False beliefs and opinions will prevent use from understanding metaphysical truths, but the error itself is not a ``metaphysical error" as it is on a totally different, and lower, plane.

The proper term, then, would be ``misconception", not ``error". However, at this time, we will leave unexplored the wisdom of counteracting one misconception with another.

\paragraph{The Swallows at Capistrano}
The swallows always seem to find their way back to Capistrano every year on March 19th. That is because they are in possession of three things: a destination, a tradition, and an initiation. They know the truth, for them, is in Capistrano. There is a tradition that reliably leads them there. And the young swallows are initiated into the intricacies of that navigation by their elders.

\paragraph{What He Said}
In his critique of Christianity as a still valid tradition\footnote{\url{https://gornahoor.net/?p=6438}}, Julius Evola brought up many good points. Just to highlight some, there is the question of a valid and continuous succession, whether there is still knowledge of the mysteries, the role of external influences on the tradition. I propose that the same standard be applied to all claimants to be the valid Tradition in the west.

We can start with the neo-pagans, the Druids, the Nietzscheans. Is there a continuous connection that goes back to an originary supernatural revelation? Do they lead a man to salvation or deliverance? Or are they a figment of someone's imagination?

Moving on, we can then ask if Evola's Buddhism meets that standard. Does it actually exist anywhere or is it one of Evola's brilliant literary creations?

Another point is the alleged egalitarianism of Christianity. Compared to what? Tantric sects are open to all, irrespective of caste or sex or nationality. But, as Emerson says, foolish consistency is the hobgoblin of little minds. Can there be a wise consistency?

\paragraph{Necker Cube}
Everyone, presumably, has seen the Necker Cube\footnote{\url{http://en.wikipedia.org/wiki/Necker_Cube}}, whose ambiguity fools the mind into projecting the cube into either the left or the right direction. Some will debate the issue; others will have a conversion experience from one viewpoint to the other. Yet others will be more conciliatory and try to see ``both sides of the issue". The scientist will describe the diagram by enumerating the line segments along with their lengths, origins, and slopes. It will be precise, but the feeling won't leave that they left something out.

So by Buddhist logic, the cube points left, it points right, it points both left and right, and it points in no direction. All points of view are relative. To transcend them is to see it sub specie aeternitatis. Either that is possible or no search for truth is possible.

\paragraph{The Search for Truth}
There can be no search for ``truth". Someone once said that America needs a good five cent cigar. Now I can embark on that search because I know what a cigar is and I know what a nickel is; I even know what a good cigar is. But how can a man seek the truth unless he knows the truth first? And if so, why would he bother to seek it?

Having a bad start, or not knowing the right direction, will likely not lead anywhere. Anyone who has ever been lost in the woods knows that a random walk bring him back to near the same spot. That is a mathematical truth. But we all know that guys hate to ask for directions.

\paragraph{False Beginnings}
In a recent interview\footnote{\url{http://sophiaperennis.org/content/response-questions-posed-n-speranskaya-journal-tradition-seyyed-hossein-nasr}}, the Traditional writer \textbf{Seyyed Hossein Nasr} said that you cannot come to Tradition by starting with \textbf{Martin Heidegger}. Yet some readers were disappointed when Gornahoor pointed that out a few weeks ago. They said we do not understand Heidegger. Our retort is that they do not understand Tradition. Or else both are true and neither is true.

\paragraph{The Method of Depth}
Those men whose primary focus is the spiritual life have more in common with each other both across nationalities and even across different traditions. Valentin Tomberg called such men Hermetists, although they may not be so explicitly or formally. That is because they are not united by their knowledge of external things, but rather by their common spiritual exercises and experiences. That is the method of depth.

\paragraph{The Case of Thomas Merton}
An example of this depth that we can point to is Thomas Merton, the American Trappist monk, who famously said he had more in common with a Zen monk than an ordinary worshiper in a pew. I haven't mentioned him much, probably because I gave away my library of Merton books to a priest over a decade ago. Merton was quite familiar with the perennial philosophy and either knew or corresponded with the second generation of Traditional writers who followed Rene Guenon and Frithjof Schuon. (I don't recall him ever mentioning Julius Evola.) The Sufi Seyyed Hossein Nasr respected Merton and mentioned him several times.

Toward the end of his life, Merton developed an interest in and appreciation for Sufism. Through his own contemplative practice and his study of Asian religions, Merton came to realize that the spiritual techniques still alive in Zen and Sufism had however been lost to Christian mysticism. Nasr describes Merton's insights:

\begin{quotex}
Since the Renaissance, much of what survived of Christian mysticism became in fact mostly individualistic, sentimental, and passive. While Merton understood the value of this type of mysticism on its own level, his gaze was set upon the great medieval Christian contemplatives whose vision was not limited in any way by the individualism that was one of the characteristics of the Renaissance. The Sufi path, in which the adept plays the active role as wayfarer upon the way that finally leads to the Beloved, like the hero in quest of the Holy Grail, while remaining passive before the grace of Heaven, was certainly attractive to Merton who was himself moving in this direction. He also thirsted for the kind of structure mystical life which the Sufi path offered in which the active and passive modes of the mystical life could be balanced on the basis of a reality that transcended the accidentality of individual existence. 

\end{quotex}
Merton's understanding of the degradation of Christian mysticism is not unlike that of Guenon or Evola. However, he did not accept that as a limitation and a reason for abandonment, as did those latter two thinkers, but rather as a motivation to learn from other traditions and incorporate what he could into his own spiritual practice. Unfortunately, Merton died before he could meet and study with Nasr personally. The task he started is languishing.

I suggest that Guenon's thinking was guided more by Sufism than by the Vedanta, if for no other reason that he understood Arabic, not Sanskrit, and actively pursued the Sufi path. His knowledge of the Vedanta was indirect and was used more for its linguistic precision. That may throw some light on the recent discussions regarding salvation and deliverance.



\flrightit{Posted on 2013-06-09 by Cologero }

\begin{center}* * *\end{center}

\begin{footnotesize}\begin{sffamily}



\texttt{Avery Morrow on 2013-06-09 at 23:31 said: }

A master (let's call him G*rd****f) sat in the forest with three students. He took out a stick and placed it in front of him, saying, ``Without calling it a stick, tell me what this is."

The first pupil said, ``You couldn't call it a slop-bucket." The master hit the student with the stick and ran away.

The second and third students were actually rocks that the pupil had mistaken for students, because it was getting very dark. Suddenly, the pupil felt alone and afraid.

In the distance, a wolf howled.


\hfill

\texttt{August on 2013-06-10 at 05:31 said: }

The wolf can smell those who make the error of imagining metaphysical errors!


\hfill

\texttt{Mihai on 2013-06-10 at 08:43 said: }

I said that in the Introduction…. he has made some unacceptable over-simplifications and generalizations regarding the West, in which he goes so far as to deny that there was any ``true intellectuality" here, minimizing the role of everything, from Christianity to Platonism.

Also, if you are referring to the predictions at the end of the book, then they have utterly failed. He explicitly said there that it is absurd to think that Easterners will succumb to the same mentality as the West. Obviously, the whole world is one big West today, and, like Saladin already told you, Guenon's view of the East had more to do with his romantic vision of it and less with the actual reality of what's been going on there. 

Like I said on the forum, ``true spirituality" is not reducible to non-dualism, nor is it a private monopoly of the orientals.


\hfill

\texttt{Jason-Adam on 2013-06-10 at 13:17 said: }

A brilliant post. Indeed, we can criticise Christianity all we want but what can we replace it with ? There is no pagan today who can show that he can trace his succession back to Antiquity which is why Evola did not practice paganism nor recommend reviving it – you simply cant revive a dead religious form, man does not make forms, God reveals them. 

I would like to learn more about Evola's personal practices, did he have a regular system of disciplines he performed every day ? I would say his Buddhism is valid, and existed or at least existed in the recent past, in Japan. While a western man can learn much from the Buddha, for myself the form is still too distant from my heritage. 

Merton interests me because in taking the best from Zen and using it to inform Catholicism, it may be the key to restoring ourselves. 

Evola destroyed Heiddegerian thinking in Ride the Tiger – it boggles me how the minds of some crazed Russians and some deranged Germanophilic Frenchmen can imagine a synthesis of Heiddeger with Evola. Whatever his errors, Evola is still so great a mind that it defames him in my view to compare him with materialistic minds like the existentialists. 

While Sufism has many profound insights, I think in the exterior world it is dangerous to endorse Islamic forms because it threatens to destroy what's left of our identity as Europeans. At least because Chinese and Japanese culture is so distant from ours, there is no risk of our becoming Japanese by adopting Zen, but Islam is very close to European tradition and thus a greater danger. Remember even Evola did say that for the exterior man, Henri Massis claim of defending the West from oriental thinking was correct.


\hfill

\texttt{Constantine Aetos on 2013-06-10 at 21:44 said: }

``At least because Chinese and Japanese culture is so distant from ours, there is no risk of our becoming Japanese by adopting Zen, "

David Carrdine did such a good job at synthesizing the European race with Asian culture. In 10 years from now, there will be more whites than Asians in Chinatown.


\hfill

\texttt{Jason-Adam on 2013-06-19 at 12:35 said: }

The problem with Christianity, by which I mean modern day practice, not the doctrine, is that there is no warrior form of the religion. In the past there was, as the writings of St Bernard clearly show – but who today in the Church defends the crusades and advocates for establishing new chivalric orders ? Because of this many warrior minds today have abandoned the Church for neopagan dead ends………..


\hfill


\end{sffamily}\end{footnotesize}
